\documentclass[11pt,a4paper]{article}
\usepackage[T1]{fontenc}
\usepackage[margin=2.5cm]{geometry}
\usepackage{amsmath,amssymb}
\usepackage{microtype}
\usepackage{helvet}
\usepackage{xcolor}
\usepackage{titlesec}
\usepackage{enumitem}
\renewcommand{\familydefault}{\sfdefault}

% Generated by tools/make_palette.py from source/palette.py.
% Do not edit by hand: edit palette.py and regenerate.
\definecolor{slideblue}{HTML}{191998}
\definecolor{slidered}{HTML}{B93333}
\definecolor{slidegray}{HTML}{686868}
\definecolor{slidegreen}{HTML}{15794B}
\definecolor{slidegold}{HTML}{A8690F}

\titleformat{\section}{\color{slideblue}\large\bfseries}{\thesection.}{.6em}{}
\titleformat{\subsection}{\color{slideblue}\normalsize\bfseries}{}{0em}{}
\titlespacing*{\subsection}{0pt}{1.5em}{.5em}
\newcommand{\slideref}[1]{\par\vspace{-.3em}{\small\itshape Supports: #1.}\par\vspace{.55em}}
\newcommand{\done}{\hfill$\blacksquare$}
% Contradiction, as introduced on the Day 1 slide. The single arrow is the
% one used for implication, so this reads as "then, and then back".
\newcommand{\contra}{\ensuremath{\rightarrow\leftarrow}}

\title{\color{slideblue}\bfseries Mathematics Bootcamp --- Proofs}
\author{Miguel V\'azquez-Carrero}
\date{Blackboard companion to the bootcamp\\[.3em]\normalsize Compilation Date: \today}

\begin{document}
\maketitle

\noindent
The slides carry the statements; these are the arguments behind them.
Everything here is meant for the board.
\vspace{1em}

\section{Day 1 --- Choice and Rationalization}

\subsection{Pairs recover revealed preference}
\slideref{``WARP and Rationalizability''}
Suppose every pair belongs to $\mathbb B$ and choice is nonempty. Under WARP,
\[
x\succeq^* y \quad\Longleftrightarrow\quad x\in C(\{x,y\}).
\]

\noindent\textbf{Proof.}
($\Leftarrow$) The pair $\{x,y\}$ is itself an observed menu containing both
alternatives, and $x$ is chosen from it. Hence $x\succeq^* y$.

\noindent
($\Rightarrow$) Let $x\succeq^* y$, so $x\in C(B)$ for some $B$ with $x,y\in B$.
Since $C(\{x,y\})\neq\varnothing$, either $x\in C(\{x,y\})$ --- and we are done
--- or $y\in C(\{x,y\})$. In the latter case WARP, applied to $B$ and
$\{x,y\}$, forces $x\in C(\{x,y\})$ as well.

\noindent
\emph{Completeness.} Every pair menu chooses something, so for any $x,y$ at
least one of $x\in C(\{x,y\})$, $y\in C(\{x,y\})$ holds; hence $x\succeq^* y$
or $y\succeq^* x$. \done

\subsection{Triples rule out cycles}
\slideref{``WARP and Rationalizability''}
Suppose the pair menus give
\[
C(\{x,y\})=\{x\},\qquad C(\{y,z\})=\{y\},\qquad C(\{x,z\})=\{z\}.
\]
No two of these menus contain the same pair, so WARP never applies across them:
with pairs alone the cycle $x$ over $y$, $y$ over $z$, $z$ over $x$ is
perfectly consistent with the axiom.

\noindent\textbf{Proof.}
Add the triple $T=\{x,y,z\}$, which a rich domain supplies. Since
$C(T)\neq\varnothing$, $T$ chooses at least one alternative.
\begin{itemize}[nosep,leftmargin=1.6em]
  \item If $x\in C(T)$: both $x$ and $z$ lie in $T$ and in $\{x,z\}$, and
  $z\in C(\{x,z\})$. WARP forces $x\in C(\{x,z\})$, contradicting
  $C(\{x,z\})=\{z\}$. \contra
  \item If $y\in C(T)$: the same argument against $C(\{x,y\})=\{x\}$. \contra
  \item If $z\in C(T)$: the same argument against $C(\{y,z\})=\{y\}$. \contra
\end{itemize}
Every case fails, so $C(T)=\varnothing$. \contra{}
The cycle is
therefore impossible, and $\succeq^*$ is transitive. \done

\subsection{WARP is equivalent to rationalizability}
\slideref{``WARP and Rationalizability''}
\textbf{Theorem.} Suppose $\mathbb B$ contains every nonempty $B\subseteq X$
with $|B|\leq 3$, and $C(B)\neq\varnothing$ for every $B\in\mathbb B$. Then $C$
is rationalizable if and only if it satisfies WARP.

\noindent\textbf{Proof.}
\emph{Necessity.} Let a rational $\succeq$ rationalize $C$. Take $x,y\in B$
with $x\in C(B)$, and $B'$ with $x,y\in B'$ and $y\in C(B')$. From $B$ we get
$x\succeq y$; from $B'$, $y\succeq z$ for every $z\in B'$. Transitivity gives
$x\succeq z$ for every $z\in B'$, so $x\in C(B')$. That is WARP.

\noindent
\emph{Sufficiency.} Let $\succeq^*$ be the revealed preference relation. By the
first lemma it is complete and satisfies $x\succeq^*y\iff x\in C(\{x,y\})$; by
the second it is transitive. So $\succeq^*$ is rational. It remains to check
$C(B)=\{x\in B:x\succeq^* y\ \text{for every}\ y\in B\}$.
\begin{itemize}[nosep,leftmargin=2.2em]
  \item[$\subseteq$] If $x\in C(B)$ then $x\succeq^* y$ for every $y\in B$,
  directly from the definition of $\succeq^*$.
  \item[$\supseteq$] Let $x\succeq^* y$ for every $y\in B$, and pick any
  $y\in C(B)$, which is nonempty. Then $x\succeq^* y$ means the pair menu
  $\{x,y\}$ chooses $x$. WARP, applied to $\{x,y\}$ and $B$ with $y\in C(B)$,
  forces $x\in C(B)$.
\end{itemize}
\done

\noindent
\emph{The key step is the last one:} if $y\in C(B)$ and $x\succeq^* y$, the
pair menu chooses $x$, and WARP carries that back into $B$.

\section{Day 3 --- Properties of Preferences}

\subsection{Monotonicity implies local nonsatiation}
\slideref{``Monotonicity and Local Nonsatiation''}
\textbf{Proposition.} On $X=\mathbb{R}_+^L$, monotonicity implies local
nonsatiation.

\noindent\textbf{Proof.}
Let $x\in\mathbb{R}_+^L$ and $\varepsilon>0$ be arbitrary. Put
$y=x+\frac{\varepsilon}{L}\mathbf 1$. Then $y\in\mathbb{R}_+^L$ and $y\gg x$,
so monotonicity gives $y\succ x$. Moreover
\[
\lVert x-y\rVert
=\sqrt{\sum_{i=1}^{L}\Bigl(\tfrac{\varepsilon}{L}\Bigr)^{2}}
=\sqrt{L\Bigl(\tfrac{\varepsilon}{L}\Bigr)^{2}}
=\frac{\varepsilon}{\sqrt L}\leq\varepsilon .
\]
So every ball $B_\varepsilon(x)$ contains a strictly better bundle.
\done

\noindent
\emph{The converse fails:} local nonsatiation does not imply monotonicity.

\subsection{Convex preferences and quasiconcave utility}
\slideref{``Convex Preferences and Quasiconcave Utility'' (Kreps 2.8b)}
Let $X$ be convex and let $u$ represent $\succeq$.

\noindent\textbf{Proof.}
For each $x\in X$,
\[
U(x)=\{y\in X:y\succeq x\}=\{y\in X:u(y)\geq u(x)\},
\]
because $u$ represents $\succeq$. So the upper contour set of $\succeq$ at $x$
\emph{is} the upper level set of $u$ at height $u(x)$. Now $\succeq$ is convex
exactly when every $U(x)$ is convex, and $u$ is quasiconcave exactly when every
upper level set is convex --- and these are the same family of sets. Replacing
the weak inequalities by strict ones throughout gives the strict version.
\done

\subsection{Lexicographic preferences have no representation}
\slideref{``Lexicographic Preferences'' (MWG Example 3.C.1)}
Let $X=\mathbb{R}_+^2$ and $x\succeq y$ iff $x_1>y_1$, or $x_1=y_1$ and
$x_2\geq y_2$.

\noindent\textbf{Proof.}
Suppose some $u$ represents $\succeq$. For each $x_1$ we have
$(x_1,2)\succ(x_1,1)$, hence $u(x_1,2)>u(x_1,1)$. By the Density Theorem
(Bartle 2.4.8) choose a rational $r(x_1)$ with
\[
u(x_1,2)>r(x_1)>u(x_1,1).
\]
The map $r$ is injective: if $x_1>x_1'$ then $(x_1,1)\succ(x_1',2)$, so
\[
r(x_1)>u(x_1,1)>u(x_1',2)>r(x_1').
\]
So $r$ is a one-to-one map from $\mathbb{R}$ into $\mathbb{Q}$. But $\mathbb{R}$
is uncountable and $\mathbb{Q}$ is countable (Bartle 1.3.6).~\contra
\done

\noindent
This is why continuity is needed: lexicographic preferences are complete,
transitive, strongly monotone and strictly convex, and still admit no
representation at all.

\subsection{A continuous representation gives continuous preferences}
\slideref{``Continuous Representations and Closed Contours''}
\textbf{Proof.}
If $f:X\to\mathbb{R}$ is continuous then $f^{-1}([c,\infty))$ and
$f^{-1}((-\infty,c])$ are closed relative to $X$ for every $c\in\mathbb{R}$.
Apply this to a continuous representation $u$ of $\succeq$:
\[
U(x)=u^{-1}\bigl([u(x),\infty)\bigr),
\qquad
L(x)=u^{-1}\bigl((-\infty,u(x)]\bigr)
\]
are then closed for every $x$. Closed contour sets are Day 3's definition of
continuity (Kreps 1.14 gives the equivalent forms), so $\succeq$ is continuous. \done

\section{Day 5 --- General Equilibrium}

\subsection{Walras' law}
\slideref{``Market Clearing''}
\textbf{Theorem.} Suppose that at prices $\{w_t,r_t\}$ the allocations
$\{c_t,a_{t+1}\}$ and $\{k_t,n_t\}$ solve the household and the firm problem,
and that in every period the labor and asset markets clear, $n_t=1$ and
$a_t=k_t$. Then the goods market clears as well, in every period.

\noindent\textbf{Proof.}
Because the allocation solves the household problem it satisfies the budget
constraint,
\[
c_t+a_{t+1}=w_t+(1+r_t)a_t .
\]
Asset-market clearing gives $a_t=k_t$ and $a_{t+1}=k_{t+1}$, so
\[
c_t+k_{t+1}=w_t+(1+r_t)k_t,
\qquad\text{that is}\qquad
c_t+k_{t+1}-(1-\delta)k_t=w_t+(r_t+\delta)k_t .
\]
The firm's first-order conditions set each factor price equal to its marginal
product. With constant returns to scale, Euler's theorem for homogeneous
functions gives
\[
w_tn_t+(r_t+\delta)k_t=A_tf(k_t,n_t),
\]
so competitive profits are zero. Using $n_t=1$,
\[
c_t+k_{t+1}-(1-\delta)k_t=w_tn_t+(r_t+\delta)k_t=A_tf(k_t,n_t),
\]
which is exactly goods-market clearing,
$c_t+k_{t+1}=A_tf(k_t,n_t)+(1-\delta)k_t$. \done

\noindent
Counting: with $T+1$ periods there are $3(T+1)$ clearing conditions but only
$2(T+1)$ prices. Walras' law resolves the apparent mismatch --- one condition
per period is redundant.

\end{document}
